\documentclass[11pt]{article}

\usepackage[utf8]{inputenc}
\usepackage[T1]{fontenc}
\usepackage[spanish,es-noshorthands]{babel}
\usepackage{amsmath,amssymb,amsfonts,mathtools}
\usepackage{bm}
\usepackage{booktabs}
\usepackage{longtable}
\usepackage{array}
\usepackage{tabularx}
\usepackage{enumitem}
\usepackage{geometry}
\usepackage{xcolor}
\usepackage{hyperref}

\geometry{letterpaper,margin=2.35cm}
\hypersetup{colorlinks=true,linkcolor=blue,citecolor=blue,urlcolor=blue}
\setlength{\emergencystretch}{2em}

\newcolumntype{L}[1]{>{\raggedright\arraybackslash}p{#1}}
\newcommand{\codefile}[1]{\begingroup\Urlmuskip=0mu plus 2mu\path{#1}\endgroup}

\newcommand{\CaputoD}{\,{}^{C}D_t^q}
\newcommand{\ii}{\mathrm{i}}
\newcommand{\R}{\mathbb{R}}
\newcommand{\sat}{\operatorname{sat}}
\newcommand{\Log}{\operatorname{Log}}
\newcommand{\sgn}{\operatorname{sgn}}
\newcommand{\diag}{\operatorname{diag}}
\newcommand{\eps}{\varepsilon}
\newcommand{\dd}{\,\mathrm{d}}

\title{Reporte técnico del repositorio\\
\large Localización y verificación de semillas candidatas para atractores ocultos en Chua fraccionario}
\author{Repositorio: \texttt{Xerkkun/Hidden-Attractors-Localization}}
\date{\today}

\begin{document}
\maketitle

\begin{abstract}
Este documento organiza, en forma defendible, la ruta matemática y numérica implementada en el repositorio \texttt{Hidden-Attractors-Localization}. El flujo central no declara ocultedad a partir de la función descriptiva; usa la función descriptiva clásica de Lur'e y la extensión tipo Machado sólo como generadores de semillas. La validación posterior se realiza mediante integración causal de Caputo con EFORK, continuación con memoria, análisis de equilibrios, criterio fraccionario de Matignon, comparación seccional, muestreo de vecindades de todos los equilibrios y mapas de cuenca. El documento separa los dos métodos de generación de semillas y describe la ruta de verificación común que decide si una semilla puede sostenerse como candidata oculta, autoexcitada o inconclusa.
\end{abstract}

\tableofcontents

\section{Alcance del reporte}

El repositorio contiene un pipeline para estudiar semillas candidatas de atractores en sistemas de Chua de orden fraccionario, principalmente en la versión no suave. La ruta implementada puede resumirse como
\[
\begin{aligned}
\text{forma Lur'e}
&\longrightarrow \text{transferencia fraccionaria}
\longrightarrow \text{Nyquist/función descriptiva}\\
&\longrightarrow \text{semilla armónica}
\longrightarrow \text{continuación}
\longrightarrow \text{simulación Caputo}\\
&\longrightarrow \text{verificación de ocultedad}.
\end{aligned}
\]

La función descriptiva, clásica o tipo Machado, no se usa como prueba de existencia de un atractor oculto. Su papel es producir condiciones iniciales informadas por balance armónico. La conclusión sobre ocultedad se posterga hasta verificar si la cuenca del atractor candidato intersecta vecindades de todos los equilibrios.

\section{Archivos principales auditados}

\begin{longtable}{@{}L{0.36\textwidth}L{0.58\textwidth}@{}}
\toprule
Archivo & Papel dentro del repositorio \\
\midrule
\endhead
\codefile{unified_nyquist_hidden_pipeline.py} & Pipeline principal. Coordina Nyquist, función descriptiva, semilla, continuación, verificación de ocultedad, cuencas, bifurcaciones, FFT/PSD y Lyapunov opcional. \\
\codefile{chua_initial_cond.py} & Núcleo matemático. Define el sistema de Chua, matrices de Lur'e, transferencia fraccionaria, funciones descriptivas, construcción de semillas y EFORK con memoria. \\
\codefile{equilibria_analysis.py} & Calcula equilibrios, jacobianos, autovalores, regiones por tramos y criterio de Matignon. \\
\codefile{harmonic_diagnostics.py} & Evalúa coeficientes de Fourier, residual armónico y razón de armónicos superiores \(\rho_H\). \\
\codefile{biased_describing_function.py} & Implementa búsqueda con función descriptiva sesgada \(N(A,\sigma_0)\), incluyendo familias clásica y Machado. \\
\codefile{multiparameter_continuation.py} & Aplica continuación por \(\eta\), posible variación de \(q\), y transporte de historia fraccionaria mediante \texttt{FractionalHistory}. \\
\codefile{run_hidden_verify_frac_hybrid.py} & Wrapper Python de verificación. Compila/ejecuta backend C, prepara configuración, lee CSV y resume la prueba de ocultedad. \\
\codefile{chua_hidden_backend.c} & Backend C pesado. Integra con EFORK, construye sección de referencia, muestrea vecindades de equilibrios y clasifica trayectorias. \\
\codefile{chua_basin_lib.c} & Backend para mapas de cuenca en cortes \(xy\), \(xz\), \(yz\) y cortes generales. \\
\codefile{seed_cloud_search.py} & Genera nubes de semillas globales fuera de vecindades de equilibrio y clasifica dinámicas no triviales. \\
\codefile{lure_candidate_manifest.py} & Organiza candidatos obtenidos por ruta clásica de Lur'e, extrae metadatos y asigna prioridad de refinamiento. \\
\codefile{lure_refined_route.py} & Ejecuta una verificación refinada de candidatos Lur'e mediante radios pequeños, direcciones propias, bolas, cascarones y mapas locales. \\
\codefile{corrida1_refined_verification.py} & Ruta refinada para candidatos Machado específicos, con etapas de tiempo, memoria y paso numérico más exigentes. \\
\bottomrule
\end{longtable}

\section{Sistema dinámico base}

El sistema principal es el Chua fraccionario no suave
\begin{align}
{}^C D_t^q x &= \alpha\bigl(y-x-f(x)\bigr),\label{eq:chua_x}\\
{}^C D_t^q y &= x-y+z,\label{eq:chua_y}\\
{}^C D_t^q z &= -\beta y-\gamma z,\label{eq:chua_z}
\end{align}
con \(0<q\leq 1\). La no linealidad usada en el modelo \texttt{piecewise} es
\begin{equation}
 f(x)=m_1x+\frac{m_0-m_1}{2}\bigl(|x+1|-|x-1|\bigr).
 \label{eq:chua_pwl}
\end{equation}
Equivalente por regiones,
\begin{equation}
 f(x)=m_1x+(m_0-m_1)\sat(x),
 \qquad
 \sat(x)=
 \begin{cases}
 -1, & x<-1,\\
 x, & |x|\leq 1,\\
 1, & x>1.
 \end{cases}
\end{equation}
Los parámetros oficiales por defecto del caso no suave son
\begin{equation}
\alpha=8.4562,
\qquad
\beta=12.0732,
\qquad
\gamma=0.0052,
\qquad
m_0=-0.1768,
\qquad
m_1=-1.1468.
\label{eq:params_official}
\end{equation}

\section{Forma tipo Lur'e}

Sea
\[
X=(x,y,z)^T,
\qquad
\sigma=r^TX=x,
\qquad
r=(1,0,0)^T.
\]
El repositorio separa la pendiente externa \(m_1\) y deja la saturación como no linealidad escalar:
\begin{equation}
{}^C D_t^q X=PX+b\psi(r^TX),
\label{eq:lure_form}
\end{equation}
con
\begin{equation}
P=
\begin{pmatrix}
-\alpha(1+m_1) & \alpha & 0\\
1&-1&1\\
0&-\beta&-\gamma
\end{pmatrix},
\qquad
b=\begin{pmatrix}-\alpha\\0\\0\end{pmatrix},
\qquad
r=\begin{pmatrix}1\\0\\0\end{pmatrix},
\label{eq:lure_matrices}
\end{equation}
y
\begin{equation}
\psi(\sigma)=A\sat(\sigma),
\qquad
A=m_0-m_1.
\label{eq:psi_sat}
\end{equation}
Esta forma es el requisito estructural para aplicar función descriptiva: el bloque lineal es tridimensional, la salida escalar es \(\sigma=x\) y la no linealidad depende sólo de \(\sigma\).

\section{Equilibrios}

Los equilibrios satisfacen
\begin{equation}
0=\alpha(y-x-f(x)),
\qquad
0=x-y+z,
\qquad
0=-\beta y-\gamma z.
\end{equation}
De las dos últimas ecuaciones,
\begin{equation}
z=y-x,
\qquad
\beta y+\gamma(y-x)=0,
\end{equation}
por tanto
\begin{equation}
y=\frac{\gamma}{\beta+\gamma}x,
\qquad
z=-\frac{\beta}{\beta+\gamma}x.
\label{eq:y_z_eq}
\end{equation}
La primera ecuación exige
\begin{equation}
y=x+f(x),
\end{equation}
de modo que
\begin{equation}
f(x^*)=-\frac{\beta}{\beta+\gamma}x^*.
\label{eq:eq_scalar}
\end{equation}
El equilibrio central es
\begin{equation}
E_0=(0,0,0).
\end{equation}
En las ramas externas \(|x^*|>1\),
\begin{equation}
f(x)=m_1x+(m_0-m_1)\sgn(x).
\end{equation}
Para la rama positiva,
\begin{equation}
m_1x_+ +(m_0-m_1)=-\frac{\beta}{\beta+\gamma}x_+,
\end{equation}
y entonces
\begin{equation}
x_+=-\frac{(m_0-m_1)(\beta+\gamma)}{(\beta+\gamma)m_1+\beta}.
\label{eq:xplus}
\end{equation}
Por simetría,
\begin{equation}
E_+=\left(x_+,\frac{\gamma}{\beta+\gamma}x_+,-\frac{\beta}{\beta+\gamma}x_+\right),
\qquad
E_-=-E_+.
\label{eq:equilibria_pm}
\end{equation}

\section{Jacobiano y criterio fraccionario de Matignon}

El jacobiano del sistema original es
\begin{equation}
J(E)=
\begin{pmatrix}
-\alpha-\alpha f'(x_E) & \alpha & 0\\
1&-1&1\\
0&-\beta&-\gamma
\end{pmatrix}.
\label{eq:jac_original}
\end{equation}
En la forma de Lur'e, esto equivale a
\begin{equation}
J(E)=P+s_E b r^T,
\label{eq:jac_lure}
\end{equation}
donde
\begin{equation}
s_E=\psi'(x_E)=
\begin{cases}
A, & |x_E|<1,\\
0, & |x_E|>1,
\end{cases}
\label{eq:slope_piecewise}
\end{equation}
excepto en las fronteras \(|x_E|=1\), donde el sistema no es diferenciable en sentido clásico.

Para el sistema lineal fraccionario conmensurado
\begin{equation}
{}^C D_t^q X=JX,
\end{equation}
el criterio de Matignon establece estabilidad local asintótica si
\begin{equation}
|\arg(\lambda_i)|>\frac{q\pi}{2},
\qquad i=1,2,3,
\label{eq:matignon}
\end{equation}
donde \(\lambda_i\) son los autovalores de \(J\). El repositorio calcula los autovalores, el margen
\begin{equation}
\Delta_i=|\arg(\lambda_i)|-\frac{q\pi}{2},
\end{equation}
y reporta
\begin{equation}
\Delta_{\min}=\min_i \Delta_i.
\end{equation}
Si \(\Delta_{\min}>0\), el equilibrio linealizado satisface Matignon. Si \(\Delta_{\min}\leq 0\), no lo satisface.

\section{Transferencia fraccionaria}

Para el bloque lineal
\begin{equation}
{}^C D_t^q X=PX+bu,
\qquad
\sigma=r^TX,
\end{equation}
la transferencia fraccionaria se escribe
\begin{equation}
\widehat W_q(\lambda)=r^T(\lambda I-P)^{-1}b,
\qquad
\lambda=s^q.
\label{eq:W_hat}
\end{equation}
El código usa la convención equivalente con signo opuesto
\begin{equation}
W_{\mathrm{code}}(\lambda)=r^T(P-\lambda I)^{-1}b=-\widehat W_q(\lambda).
\label{eq:W_code}
\end{equation}
Sobre el eje de frecuencia fraccionario,
\begin{equation}
\lambda=(j\omega)^q=\omega^q\exp\left(j\frac{q\pi}{2}\right),
\qquad \omega>0.
\label{eq:lambda_frequency}
\end{equation}
Por tanto,
\begin{equation}
W_{\mathrm{code}}(j\omega)=r^T\left(P-(j\omega)^qI\right)^{-1}b.
\label{eq:W_code_freq}
\end{equation}
La variable \(j\omega\) no se introduce como en sistemas enteros; siempre se evalúa mediante \((j\omega)^q\).

\section{Puente Weyl--Caputo}

La función descriptiva, Nyquist y el balance armónico presuponen señales sinusoidales persistentes. En sistemas autónomos de Caputo no se debe concluir la existencia de ciclos periódicos exactos no constantes sólo por obtener una solución armónica formal. La interpretación usada por el repositorio es la siguiente:
\begin{enumerate}[label=\alph*)]
\item Caputo describe el sistema causal con memoria desde un tiempo inicial \(t_0\).
\item Weyl/Liouville--Weyl permite representar regímenes periódicos estacionarios ideales en frecuencia.
\item La diferencia entre ambos marcos aparece como término de memoria o discrepancia transitoria.
\item El balance armónico en Weyl genera semillas; la validación se hace después en Caputo mediante integración causal.
\end{enumerate}
Así, las semillas armónicas se interpretan como aproximaciones estacionarias o asintóticas, no como pruebas exactas de ciclos límite.

\section{Método I: semilla por función descriptiva clásica}

\subsection{Función descriptiva de la saturación}

Se toma una entrada armónica
\begin{equation}
\sigma(t)=a\cos\theta,
\qquad \theta=\omega t.
\end{equation}
La salida no lineal es
\begin{equation}
y(\theta)=\psi(a\cos\theta).
\end{equation}
El primer armónico se escribe
\begin{equation}
y_1(\theta)=Y_1\cos\theta+\widetilde Y_1\sin\theta.
\end{equation}
Como \(\psi\) es impar y sin memoria, la función descriptiva clásica es real:
\begin{equation}
N(a)=\frac{Y_1}{a}.
\end{equation}
Para \(\psi(\sigma)=A\sat(\sigma)\), se obtiene
\begin{equation}
N(a)=
\begin{cases}
A, & 0<a\leq 1,\\[6pt]
\displaystyle
\frac{2A}{\pi}\left[\arcsin\left(\frac{1}{a}\right)+\frac{\sqrt{a^2-1}}{a^2}\right], & a>1.
\end{cases}
\label{eq:N_classic_sat}
\end{equation}

\subsection{Condición de Nyquist}

La condición armónica es
\begin{equation}
1+W_{\mathrm{code}}(j\omega)N(a)=0.
\label{eq:harmonic_classic}
\end{equation}
Como \(N(a)\in\R\), se impone
\begin{equation}
\operatorname{Im} W_{\mathrm{code}}(j\omega_0)=0,
\label{eq:im_zero}
\end{equation}
\begin{equation}
N(a_0)=k_0,
\qquad
k_0=-\frac{1}{\operatorname{Re}W_{\mathrm{code}}(j\omega_0)}.
\label{eq:k0_classic}
\end{equation}
El código escanea \(\omega\in[10^{-4},10]\), localiza cambios de signo de \(\operatorname{Im}W\), refina raíces con un método tipo Brent y ordena las ramas por \(k\). Sólo se aceptan ramas compatibles:
\begin{equation}
0<k_0\leq A.
\end{equation}

\subsection{Construcción de la semilla clásica}

Con \(k_0=N(a_0)\), se define
\begin{equation}
P_0=P+k_0br^T.
\label{eq:P0_classic}
\end{equation}
Sea
\begin{equation}
\lambda_0=(j\omega_0)^q.
\end{equation}
El autovector complejo \(v\) satisface aproximadamente
\begin{equation}
(P_0-\lambda_0I)v=0,
\qquad
r^Tv=1.
\label{eq:eig_seed}
\end{equation}
La semilla con fase \(\theta\) es
\begin{equation}
X_{\mathrm{seed}}(\theta)=a_0\operatorname{Re}\left(v e^{j\theta}\right).
\label{eq:seed_classic}
\end{equation}
En la ruta clásica base se toma \(\theta=0\); en rutas de exploración se barren varias fases.

\section{Método II: extensión tipo Machado}

\subsection{Función descriptiva fraccionaria auxiliar}

El repositorio implementa una extensión inspirada en la propuesta de Machado. Si \(N(a)\) es la función descriptiva clásica, se define
\begin{equation}
N_\mu(a)=N(a)^\mu,
\qquad \mu>0,
\label{eq:machado_real}
\end{equation}
para el caso real positivo de la saturación no suave. La condición armónica queda
\begin{equation}
1+W_{\mathrm{code}}(j\omega)N_\mu(a)=0.
\label{eq:machado_closure}
\end{equation}
Como \(N_\mu(a)\in\R\), la frecuencia sigue saliendo de
\begin{equation}
\operatorname{Im} W_{\mathrm{code}}(j\omega_j)=0,
\end{equation}
y la ganancia equivalente de la rama es
\begin{equation}
k_j=-\frac{1}{\operatorname{Re}W_{\mathrm{code}}(j\omega_j)}.
\end{equation}
La amplitud Machado se obtiene resolviendo
\begin{equation}
N(a_{\mu,j})^\mu=k_j,
\end{equation}
o equivalentemente
\begin{equation}
N(a_{\mu,j})=k_j^{1/\mu}.
\label{eq:machado_amp}
\end{equation}
La condición de compatibilidad es
\begin{equation}
0<k_j\leq A^\mu.
\label{eq:machado_compatibility}
\end{equation}
Si esta desigualdad falla, el par \((j,\mu)\) se descarta.

\subsection{Potencia compleja para casos sesgados}

Para búsquedas sesgadas o funciones descriptivas complejas, el repositorio usa
\begin{equation}
N_\mu=\exp\{\mu\Log_\ell(N)\},
\label{eq:machado_complex}
\end{equation}
donde \(\ell\) indica la rama logarítmica. En rama principal,
\begin{equation}
\Log_0(N)=\ln|N|+j\arg(N),
\qquad \arg(N)\in(-\pi,\pi].
\end{equation}
Esta versión se usa para familias \texttt{machado\_biased} y permite no linealidades donde la componente fundamental tenga desfase.

\subsection{Semilla Machado}

Para cada rama \(j\), orden \(\mu\) y fase \(\theta\), la semilla es
\begin{equation}
X_{\mu,j,\theta}=a_{\mu,j}\operatorname{Re}\left(v_j e^{j\theta}\right),
\label{eq:seed_machado}
\end{equation}
donde \(v_j\) es el autovector asociado a \(\lambda_j=(j\omega_j)^q\), normalizado por \(r^Tv_j=1\). La diferencia frente al método clásico no está en la verificación, sino en la amplitud y en la exploración de \(\mu\) y \(\theta\).

\section{Función descriptiva sesgada}

El flujo extendido considera
\begin{equation}
\sigma(t)=\sigma_0+A\cos(\omega t).
\end{equation}
Se calculan los coeficientes de Fourier de
\begin{equation}
y(\theta)=\psi(\sigma_0+A\cos\theta),
\end{equation}
con convención
\begin{equation}
y(\theta)=y_0+\sum_{k\geq 1}\left(a_k\cos(k\theta)+b_k\sin(k\theta)\right),
\end{equation}
\begin{equation}
Y_k=a_k-jb_k.
\end{equation}
La función descriptiva sesgada es
\begin{equation}
N(A,\sigma_0)=\frac{Y_1}{A}.
\label{eq:N_biased}
\end{equation}
Para reconstruir la semilla, se separa una parte constante y una parte armónica. La parte media \(\bar X\) se calcula con
\begin{equation}
P\bar X+b\bar y=0,
\qquad
r^T\bar X=\sigma_0.
\label{eq:dc_biased}
\end{equation}
La componente fundamental \(V\) satisface
\begin{equation}
(\lambda I-P)V=bY_1,
\qquad
r^TV=A,
\qquad
\lambda=(j\omega)^q.
\label{eq:harm_biased}
\end{equation}
La semilla reconstruida es
\begin{equation}
X_{\mathrm{seed}}=\bar X+\operatorname{Re}(Ve^{j\theta}).
\label{eq:seed_biased}
\end{equation}

\section{Diagnóstico de armónicos superiores}

Para evaluar si la hipótesis de primer armónico es razonable, el repositorio calcula
\begin{equation}
R=1+W_{\mathrm{code}}(j\omega)N_{\mathrm{eff}},
\label{eq:residual_general}
\end{equation}
donde
\begin{equation}
N_{\mathrm{eff}}=
\begin{cases}
N, & \text{familia clásica},\\
\exp\{\mu\Log(N)\}, & \text{familia Machado}.
\end{cases}
\end{equation}
Además, calcula la razón
\begin{equation}
\rho_H=
\frac{\displaystyle\sum_{k=2}^{K}|W_{\mathrm{code}}(jk\omega)|\,|Y_k|}
{|W_{\mathrm{code}}(j\omega)|\,|Y_1|+\delta},
\label{eq:rhoH}
\end{equation}
con \(\delta>0\) pequeño para evitar división por cero. Si \(\rho_H\ll1\), la componente fundamental domina después del bloque lineal. Si \(\rho_H\) es grande, la semilla puede explorarse, pero la predicción armónica es débil.

\section{Continuación numérica}

\subsection{Continuación clásica en \texorpdfstring{\(\eps\)}{epsilon}}

La familia de continuación clásica es
\begin{equation}
{}^C D_t^q X=P_0X+\eps b\delta(r^TX),
\qquad
\delta(\sigma)=\psi(\sigma)-k\sigma,
\qquad
\eps\in[0,1].
\label{eq:continuation_eps}
\end{equation}
Para \(\eps=0\),
\begin{equation}
{}^C D_t^qX=P_0X,
\end{equation}
que corresponde al sistema lineal auxiliar. Para \(\eps=1\),
\begin{align}
P_0X+b\delta(r^TX)
&=P X+kbr^TX+b\psi(r^TX)-kb r^TX\\
&=PX+b\psi(r^TX),
\end{align}
por lo que se recupera el sistema original.

\subsection{Continuación multiparamétrica}

El flujo extendido usa una homotopía adicional en la no linealidad no suave:
\begin{equation}
\psi_\eta(\sigma)=A\left[(1-\eta)\tanh\left(\frac{\sigma}{w}\right)+\eta\sat(\sigma)\right],
\qquad
\eta\in[0,1].
\label{eq:eta_homotopy}
\end{equation}
Así, \(\eta=0\) representa un sistema suavizado y \(\eta=1\) representa el sistema no suave original. La agenda de continuación puede combinar etapas en \(\eta\), \(q\) y un parámetro auxiliar \(\chi\), que puede ser \(\mu\) o \(\sigma_0\), según la familia de semilla.

\section{Integración EFORK y memoria}

El sistema de Caputo se integra mediante EFORK con memoria truncada. Para una componente escalar, el término de memoria discreto tiene la forma
\begin{equation}
M_n\approx
\frac{1}{h\Gamma(2-q)}
\sum_{j=\max(0,n-\nu)}^{n-1}
(x_{j+1}-x_j)
\left[(t_n-t_j)^{1-q}-(t_n-t_{j+1})^{1-q}\right],
\label{eq:memory_term}
\end{equation}
con
\begin{equation}
\nu=\left\lceil\frac{L_m}{h}\right\rceil.
\end{equation}
Esto implica que el estado fraccionario no queda determinado por \(X(t_n)\) únicamente. En una continuación coherente se debe transportar la ventana
\begin{equation}
\{X(t_{n-\nu}),X(t_{n-\nu+1}),\ldots,X(t_n)\}.
\label{eq:memory_window}
\end{equation}
El repositorio distingue dos modos:
\begin{description}[leftmargin=3.2cm]
\item[\texttt{point}] transporta sólo el último punto. Es barato, pero reinicia artificialmente la memoria.
\item[\texttt{window}] transporta una ventana de historia de longitud aproximada \(L_m\). Es el modo defendible para Caputo con memoria truncada.
\end{description}

\section{Ruta común de verificación de semillas}

La verificación final es la misma para semillas clásicas, Machado, sesgadas o de nube. El procedimiento es:
\begin{enumerate}[label=\arabic*.]
\item Calcular todos los equilibrios del sistema.
\item Clasificarlos localmente mediante el criterio de Matignon.
\item Integrar la semilla candidata en el sistema original de Caputo.
\item Construir una referencia geométrica del atractor candidato usando una sección.
\item Muestrear vecindades de cada equilibrio.
\item Integrar todas las condiciones iniciales de prueba.
\item Clasificar cada trayectoria.
\item Decidir si la cuenca candidata intersecta o no las vecindades de equilibrio muestreadas.
\end{enumerate}

\subsection{Sección de referencia}

La referencia del atractor candidato se construye sobre la sección
\begin{equation}
x=0,
\qquad
\dot x>0.
\end{equation}
Numéricamente se detectan cruces
\begin{equation}
x_{k-1}<0,
\qquad
x_k\geq 0,
\end{equation}
y se interpola el punto \((y,z)\). Después de descartar el transitorio, se obtiene una nube
\begin{equation}
\mathcal S_*=\{(y_i,z_i)\}_{i=1}^{N_s}.
\end{equation}
Esta nube no es un mapa de retorno exacto; es una huella geométrica operacional del atractor candidato.

\subsection{Muestreo alrededor de equilibrios}

Para cada equilibrio \(E_j\), se generan condiciones iniciales en bolas
\begin{equation}
X_0\in B(E_j,r),
\end{equation}
para varios radios \(r\). También existen rutas refinadas que usan cascarones, direcciones coordenadas, diagonales y direcciones propias del jacobiano.

\subsection{Clasificación de trayectorias}

Cada trayectoria se clasifica como:
\begin{description}[leftmargin=3.0cm]
\item[\texttt{EQ}] queda dentro del criterio operacional de cercanía a algún equilibrio.
\item[\texttt{DIV}] excede un radio de divergencia \(R_{\mathrm{DIV}}\).
\item[\texttt{TARGET}] su sección se parece a la referencia \(\mathcal S_*\).
\item[\texttt{OTHER}] es acotada y no coincide con la referencia.
\item[\texttt{UNKNOWN}] no hay puntos de sección suficientes o la evidencia es insuficiente.
\end{description}

La condición \texttt{TARGET} se decide mediante
\begin{equation}
\frac{\#\{p\in\mathcal S_{\mathrm{test}}:\operatorname{dist}(p,\mathcal S_*)\leq \mathrm{SEC\_TOL}\}}
{\#\mathcal S_{\mathrm{test}}}
\geq \mathrm{HIT\_FRAC\_REQ}.
\label{eq:target_fraction}
\end{equation}

\section{Criterio operativo de ocultedad}

La definición operativa es
\begin{equation}
\mathcal A \text{ es oculto}
\quad\Longleftrightarrow\quad
\mathcal B(\mathcal A)\cap U(E_j)=\varnothing
\quad\text{para toda vecindad abierta }U(E_j)
\end{equation}
de cada equilibrio \(E_j\). Numéricamente no se pueden probar todas las vecindades abiertas; el repositorio prueba una familia finita de radios. La regla conservadora es:
\begin{equation}
\text{si existe }E_j,r\text{ con }\texttt{TARGET}>0,
\text{ no declarar ocultedad fuerte.}
\end{equation}
La decisión práctica es:
\begin{itemize}
\item \texttt{hidden-like}: semilla acotada, no trivial y no cercana a equilibrios, pero sin verificación completa.
\item \texttt{numerically supported hidden}: no hay \texttt{TARGET} desde las vecindades muestreadas de ningún equilibrio.
\item \texttt{self-excited}: hay intersección robusta con vecindades de algún equilibrio.
\item \texttt{inconclusive}: hay ambigüedad por tiempos, radios, memoria, paso, sección o umbrales.
\end{itemize}

\section{Diferencia entre ambos métodos}

La diferencia entre la ruta clásica y la ruta Machado está en la generación de semillas, no en la verificación final.

\begin{center}
\small
\begin{tabular}{@{}L{0.26\textwidth}L{0.31\textwidth}L{0.34\textwidth}@{}}
\toprule
Aspecto & Método clásico de Lur'e & Método tipo Machado \\
\midrule
Función descriptiva & \(N(a)\) & \(N_\mu(a)=N(a)^\mu\) o \(\exp\{\mu\Log(N)\}\) \\
Parámetro auxiliar & Ninguno, salvo rama y fase & Orden descriptivo \(\mu>0\), rama logarítmica y fase \\
Condición de cierre & \(1+W_qN(a)=0\) & \(1+W_qN_\mu(a)=0\) \\
Amplitud & \(N(a_0)=k\) & \(N(a_\mu)=k^{1/\mu}\) \\
Exploración & Pocas ramas Nyquist & Ramas \(j\), malla de \(\mu\), fases \(\theta\), y opción sesgada \\
Papel matemático & Semilla armónica clásica & Ampliación heurística de semillas \\
Verificación & Integración Caputo, cuencas y vecindades de equilibrios & La misma verificación que la ruta clásica \\
Conclusión permitida & No prueba ocultedad por sí sola & No prueba ocultedad por sí sola \\
\bottomrule
\end{tabular}
\end{center}

\section{Rutas refinadas}

\subsection{Ruta refinada Lur'e}

La ruta \texttt{lure\_candidate\_manifest.py} recolecta candidatos clásicos, extrae
\begin{equation}
q,\omega,k,A,\theta,X_{\mathrm{seed}},X_{\mathrm{final}},
\end{equation}
y resume los contactos con equilibrios:
\begin{equation}
\texttt{target\_from\_Eminus},\qquad
\texttt{target\_from\_E0},\qquad
\texttt{target\_from\_Eplus}.
\end{equation}
Después asigna prioridad:
\begin{equation}
\texttt{contact\_strong},\quad
\texttt{contact\_weak},\quad
\texttt{no\_target\_under\_sample},\quad
\texttt{insufficient\_data}.
\end{equation}
La ruta \texttt{lure\_refined\_route.py} vuelve a muestrear vecindades usando radios pequeños, bolas, cascarones y direcciones propias del jacobiano. Esto fortalece o bloquea la afirmación de ocultedad.

\subsection{Ruta refinada Machado}

La ruta Machado usa barridos \texttt{machado\_sweep} y \texttt{machado\_sweep\_fast}. Para candidatos seleccionados, \texttt{corrida1\_refined\_verification.py} aplica etapas con diferentes tiempos finales, pasos y longitudes de memoria. El objetivo es verificar si un candidato Machado sobrevive cuando se incrementa el rigor numérico.

\section{Diagnósticos auxiliares}

El repositorio produce diagnósticos que ayudan a caracterizar la dinámica, pero no sustituyen la prueba de ocultedad:
\begin{itemize}
\item FFT directa para identificar frecuencias dominantes.
\item PSD de Welch opcional.
\item Exportación a TISEAN opcional.
\item Estimación operacional de exponentes de Lyapunov con backend C.
\item Diagramas de bifurcación en parámetros y en \(q\).
\item Mapas de cuenca en cortes \(xy\), \(xz\), \(yz\).
\end{itemize}
Una semilla con exponente de Lyapunov positivo puede indicar caos numérico, pero eso no decide si el atractor es oculto. La ocultedad depende de la cuenca respecto a los equilibrios.

\section{Advertencias técnicas}

\subsection{La clase \texttt{EQ} es operacional}

En algunas pruebas, una trayectoria puede clasificarse como \texttt{EQ} si permanece dentro de una bola de radio \(\mathrm{EPS\_EQ}\) durante \(\mathrm{CAP\_WIN}\) pasos. Si el equilibrio es inestable, una condición inicial extremadamente cercana puede permanecer ahí un tiempo finito antes de escapar. Por tanto, \texttt{EQ} debe leerse como captura operacional finita, no necesariamente como convergencia asintótica probada. Una versión más estricta debe verificar persistencia terminal después del transitorio.

\subsection{La memoria está truncada}

El integrador usa \(L_m<\infty\). Esto permite cómputo viable, pero introduce una aproximación de memoria corta respecto a Caputo con historia completa. Al defender resultados, se debe reportar \(h\), \(L_m\), transitorio, tiempo final y número de puntos de memoria.

\subsection{La función descriptiva es heurística}

Ni \(N(a)\) ni \(N_\mu(a)\) prueban la existencia de un atractor. Ambas generan semillas. La conclusión científica proviene de la integración y de las cuencas.

\section{Conclusiones defendibles}

\begin{enumerate}[label=\arabic*.]
\item El repositorio está estructurado correctamente alrededor de una forma tipo Lur'e compatible con función descriptiva: \({}^C D_t^qX=PX+b\psi(r^TX)\).
\item La transferencia fraccionaria se evalúa con \(\lambda=(j\omega)^q\), no con \(j\omega\) entero.
\item La ruta clásica usa \(N(a)\) para obtener \((\omega_0,k_0,a_0)\) y construir una semilla armónica.
\item La ruta Machado usa \(N_\mu(a)=N(a)^\mu\) o \(\exp\{\mu\Log N\}\) para ampliar la exploración de amplitudes y fases.
\item En ambos métodos, la función descriptiva sólo produce semillas; no prueba ocultedad.
\item La continuación con ventana de memoria es más defendible que la continuación punto a punto para Caputo.
\item La verificación de ocultedad exige simular desde vecindades de todos los equilibrios y comprobar ausencia de trayectorias \texttt{TARGET} hacia el atractor candidato.
\item Si aparece \texttt{TARGET} desde una vecindad de equilibrio, el candidato no debe declararse oculto bajo esa configuración.
\item FFT, PSD, bifurcaciones y Lyapunov caracterizan la dinámica, pero no sustituyen la prueba de cuenca.
\item La conclusión final siempre debe reportarse con \(q,h,L_m,T_{\mathrm{burn}},T_{\mathrm{final}}\), radios, número de muestras, criterio seccional y umbrales.
\end{enumerate}

\section{Checklist para defender una semilla candidata}

Antes de afirmar que una semilla sostiene un atractor oculto, debe poder llenarse la siguiente lista:
\begin{enumerate}[label=\arabic*.]
\item Sistema y parámetros oficiales fijados.
\item Forma de Lur'e escrita explícitamente.
\item Equilibrios \(E_0,E_+,E_-\) calculados.
\item Jacobiano y Matignon reportados para cada equilibrio.
\item Transferencia \(W_q\) evaluada con \(\lambda=(j\omega)^q\).
\item Método de semilla identificado: clásico, Machado, sesgado o nube.
\item Residual armónico \(|1+W_qN_{\mathrm{eff}}|\) reportado.
\item Diagnóstico \(\rho_H\) reportado si se usa balance armónico.
\item Semilla inicial y semilla final después de continuación reportadas.
\item Integrador, paso, memoria y tiempos reportados.
\item Referencia seccional construida y número de puntos reportado.
\item Radios y modos de muestreo alrededor de equilibrios reportados.
\item Tabla por equilibrio y radio con \texttt{EQ}, \texttt{DIV}, \texttt{TARGET}, \texttt{OTHER}, \texttt{UNKNOWN}.
\item Criterio explícito para aceptar o rechazar ocultedad.
\item Cuencas y trayectorias graficadas.
\item Resultado final etiquetado de forma conservadora: oculto soportado numéricamente, autoexcitado o inconcluso.
\end{enumerate}

\end{document}
